\documentclass{beamer}

\usepackage[utf8]{inputenc}
\usepackage[english]{babel}
\usepackage{textpos}
\usepackage{graphicx}
\usepackage{textcomp}
\usepackage{animate}
\usepackage{comment}
%\usepackage{multicol}
%\setlength{\columnsep}{5cm}
\usepackage{amsmath}
\usepackage{amsfonts}
\usepackage{verbatim}
\usepackage{varwidth}
\usepackage{caption}
\usepackage{dcolumn}
\usepackage{amssymb}
\usepackage{algorithm}
\usepackage{algorithmic}
%\usepackage{physics}
\usepackage{hyperref}
\usepackage{scalerel,stackengine}
%\usepackage{fancyvrb}


\newcommand{\code}[1]{\texttt{#1}}
\newcommand{\shp}{$>>>$ }
\newcommand{\tab}{\hspace*{22pt}}

\usetheme{Madrid}

\title{NumPy and Matplotlib}
\author{EEC 3150}
\date{}

\institute{Trevecca Nazarene University}
 


%%%%%%%%% BEGIN DOCUMENT %%%%%%%%%
%%%%%%%%% BEGIN DOCUMENT %%%%%%%%%
%%%%%%%%% BEGIN DOCUMENT %%%%%%%%%

\begin{document}

% get title page 
\frame{\titlepage}
 
%%%%%%%%%%
%%%%%%%%%%
\begin{frame}	
	\frametitle{Outline}
	\tableofcontents	
\end{frame}


%%%%%%%%%%
%%%%%%%%%%
\section{NumPy}

\begin{frame}	
	\center{\Huge{NumPy: \\ Numerical Python}}
\end{frame}

\subsection{Intro to NumPy}

\begin{frame}	
	\frametitle{Intro to NumPy}
	
	\begin{block}{NumPy:}
		NumPy is one of the most commonly used third-party libraries in Python \\
		\href{https://numpy.org/doc/stable/}{https://numpy.org/doc/stable/}
		
		\vspace{10pt}
		Key features:
		\begin{itemize}
			\item Extensive library of tools for mathematical computations
			\begin{itemize}
				\item Trigonometric, exponential, and special functions
				\item Statistics and random number generation
				\item Linear algebra and Fourier analysis (and more)
			\end{itemize}
			\item Bread and butter: \code{ndarray} objects
			\begin{itemize}
				\item $n$-dimensional arrays
				\item Lists but infinitely better (how you wish lists worked)
			\end{itemize}
		\end{itemize}
		
		\vspace{10pt}
		Importing numpy: \\
		\code{import numpy as np}
	\end{block}
	
\end{frame}


%%%%%%%%%%
%%%%%%%%%%
\subsection{Creating \code{ndarray}'s}

\begin{frame}	
	\frametitle{Creating \code{ndarray}'s}
	
	\begin{block}{Functions for creating arrays:}
		\begin{scriptsize}
		
%		NumPy provides many functions for creating arrays of different kinds
		\begin{itemize}
			\item \code{np.array(x)} will convert the list \code{x} into an array \\
					The function accepts an additional keyword argument \code{dtype} that allows for the specification of data type the array should have (ex: np.int32, np.float64)
			\item \code{np.zeros(n)} will create an array of length $n$ (int) of all zeros \\
					You can also pass a tuple instead of an integer to create a multi-dimensional array (ex: \code{np.zeros((m,n))} creates and $m\times n$ array)
			\item \code{np.ones(n)} (like above, but all ones)
			\item \code{np.full(n,value)} (like above, but fills array with \code{value})
			\item \code{eye(n)} creates an $n\times n$ identity matrix
			\item \code{np.arange(a,b,step)} same as Python \code{range} but allows floats
			\item \code{np.linspace(a,b,n)} creates an array of \code{n} evenly-spaced points from \code{a} to \code{b}
			\item \code{np.random.rand(n)} creates a length \code{n} array of random numbers b/t 0 and 1 \\
					This uses the \code{np.random} library (more on that later)
		\end{itemize}
		Some additional points:
		\begin{itemize}
			\item The different dimensions of a numpy array are called \emph{axes} (\code{axis} is often a keyword argument to many functions)
			\item Use the \code{shape} attribute to get the length of an array along each axis (ex: \code{x.shape})
		\end{itemize}
		
%		Equivalent: \code{2*np.ones((3,3))} and \code{np.full((3,3),2)}
		\end{scriptsize}
	\end{block}
	
\end{frame}

\begin{frame}	
	\frametitle{Creating \code{ndarray}'s}
	
	\begin{exampleblock}{2 ways of creating a $3\times3$ array of all 5's:}
		\begin{itemize}
			\item np.full((3,3),5)
			\item 5*np.ones((3,3))
		\end{itemize}
	\end{exampleblock}
	
	\begin{exampleblock}{Row vs. column vs. 1-D:}
		These are all different objects:
		\begin{itemize}
			\item np.ones(5) \\
					This is a 1-D array
			\item np.ones((5,1)) \\
					This is a 2-D array (column)
			\item np.ones((1,5)) \\
					This is a 2-D array (row)
		\end{itemize}
	\end{exampleblock}
	
\end{frame}

%%%%%%%%%%
%%%%%%%%%%
\subsection{Math operations on \code{ndarray}'s}

\begin{frame}	
	\frametitle{Math operations \code{ndarray}'s}
	
	\begin{block}{Math operations \code{ndarray}'s:}
%		\begin{footnotesize}
		
		Numpy arrays can handle mathematical operations in an intuitive, element-wise way:
		\begin{itemize}
			\item We can perform virtually any mathematical operation on any combination of arrays and scalars (as long as their lengths in each dimension match)
			\item The exception to this is Boolean operations, but Numpy has functions for these anyway
			\item We can also apply virtually any of Numpy's mathematical functions on arrays in an element-wise fashion
			\item Numpy can even add arrays of different sizes in certain contexts
		\end{itemize}
%		\end{footnotesize}
	\end{block}
	
\end{frame}

\begin{frame}	
	\frametitle{Math operations \code{ndarray}'s}
	
	\begin{exampleblock}{Plot some random numbers:}
		\begin{itemize}
			\item Create an array of random numbers b/t $-a$ and $a$
			\item Create and plot a linearly-spaced array of numbers with the above random numbers added as noise
		\end{itemize}
	\end{exampleblock}
	
	\begin{exampleblock}{Plot sin and cos:}
		Plot sin$(x)$ and cos$(x)$ from 0 to $2\pi$
	\end{exampleblock}
	
\end{frame}

% row vs. column vs. 1-d

%%%%%%%%%%
%%%%%%%%%%
\subsection{Indexing \code{ndarray}'s}

\begin{frame}	
	\frametitle{Indexing \code{ndarray}'s}
	
	\begin{block}{Indexing \code{ndarray}'s}
		\begin{scriptsize}
		Standard list indexing:
		\begin{itemize}
			\item Index an element: \code{x[1]}
			\item Get a slice: \code{x[4:12]}
			\item Index from the end: \code{x[-3]}
		\end{itemize}
		
		Numpy also supports \emph{fancy indexing}:
		\begin{itemize}
			\item Indexing with \code{np.newaxis} will add another axis to the array
			\item Index a 2-D array: \code{x[2,5]}
			\item Index an entire row, column respectively: \code{x[2,:]}, \code{x[:,5]}
			\item Index a list of specific elements: \code{x[[1,4,8]]}
			\item Index all elements with values greater than 0: \code{x[x>0]} \\
			\begin{itemize}
				\scriptsize
				\item How it works: \code{x>0} creates a boolean mask used to index the \code{x} array
				\item This works for multi-dimensional arrays too
				\item Limitation: you can't do multiple conditionals in the same fancy index
				% (but you can still accomplish the same functionality using Numpy's functions for logical operations)
			\end{itemize}
			\item Get the indices (not values) where an array is greater than 0: \code{np.where(x>0)}
		\end{itemize}
		
		In addition to accessing array elements with the above techniques, you can also SET elements with the above techniques
		\end{scriptsize}
	\end{block}
	
\end{frame}

\begin{frame}	
	\frametitle{Indexing \code{ndarray}'s}
	
	\begin{block}{Conditional masking vs. np.where()}
		Let \code{x} be a numpy array:
		\begin{itemize}
			\item \code{x[x>0]} returns the VALUES where \code{x} is greater than 0
			\item \code{np.where(x>0)} returns the INDICES where \code{x} is greater than 0
		\end{itemize}
	\end{block}
	
	\begin{block}{np.where() on multi-dimensional arrays}
		\begin{itemize}
			\item Running a \code{np.where()} on a 2D (or higher) array will produce a tuple of index arrays
			\item each tuple corresponds to the indices of the respective axis
			\item Indexing the original array with this tuple of index arrays will produce a 1D array
		\end{itemize}
	\end{block}
	
\end{frame}

\begin{frame}	
	\frametitle{Indexing \code{ndarray}'s}
	
	\begin{exampleblock}{Set negative elements to 0:}
		\begin{itemize}
			\footnotesize
			\item Generate some random data
			\item Set all negative values to 0
		\end{itemize}
	\end{exampleblock}
	
	\begin{exampleblock}{Index with multiple conditions:}
		\begin{itemize}
			\footnotesize
			\item Generate some random data
			\item Using several boolean masks, index the values such that $a<x<b$
			\item There are several ways to do this
			\begin{itemize}
				\footnotesize
				\item Use numpy's \code{np.logical\_and(x,y)} function
				\item Repeatedly index and overwrite the same array with different boolean masks
			\end{itemize}
		\end{itemize}
	\end{exampleblock}
	
	\begin{exampleblock}{Plot positive sin$(x)$:}
		\begin{itemize}
			\footnotesize
			\item Create \code{x,y} arrays for plotting sin$(x)$ from 0,$4\pi$
			\item Only plot values where sin$(x)\ge0$
			\item Plot with dots
		\end{itemize}
	\end{exampleblock}
	
\end{frame}

% some helpful array functions
% reshape, flatten, unique, sort, argsort, transpose

% math functions
% sin, cos, tan, exp, inverse

% set functions
% union, intersection

% boolean functions
% logical\_and

% rng, np.random
% rand, randint, uniform, normal, multivariate\_normal

% statistics functions
% mean, median, mode, std, var, cov, corrcoef, histogram

% linear algebra, np.linalg
% eye, dot, det, inv





\end{document}




